% from ega3-1-fr.tex
\oldpage[III]{97}
\begin{corollary}[2.1.4]
\label{III.2.1.4-fr}
Si $X$ est un schéma quasi-compact et si les $U_i=X_{f_i}$ sont affines,
il existe un isomorphisme canonique fonctoriel en $\mathcal F$, de degré $0$,
\begin{equation}
\label{III.2.1.4.1-fr}
\tag{2.1.4.1}
H^p(U,\mathcal F(*))\xrightarrow{\sim}H^{p+1}((\mathbf f),M)
\qquad\text{pour }p\geq1,
\end{equation}
et une suite exacte fonctorielle en $\mathcal F$
\begin{equation}
\label{III.2.1.4.2-fr}
\tag{2.1.4.2}
0\longrightarrow H^0((\mathbf f),M)\longrightarrow M
\longrightarrow H^0(U,\mathcal F(*))
\longrightarrow H^1((\mathbf f),M)\longrightarrow0,
\end{equation}
où tous les homomorphismes sont de degré $0$.
\end{corollary}

\begin{proof}
Il suffit en effet d'appliquer (1.4.1) au résultat de
(\hyperref[III.2.1.3-fr]{2.1.3}).
\end{proof}

La proposition «~locale~» analogue à
(\hyperref[III.2.1.3-fr]{2.1.3}) est la suivante~:

\begin{proposition}[2.1.5]
\label{III.2.1.5-fr}
Soient $S$ un anneau gradué à degrés positifs,
$f_i$ $(1\leq i\leq r)$ un élément homogène de $S_+$ de degré $d_i$,
$M$ un $S$-module gradué. Soit $X=\operatorname{Proj}(S)$ le spectre
premier homogène de $S$ et posons $U_i=D_+(f_i)$,
$U=\bigcup_iU_i$,
\[
H^\bullet(U,\widetilde M(*))
=\bigoplus_{n\in\mathbb Z}H^\bullet(U,(M(n))^\sim).
\]
Il existe alors des isomorphismes canoniques fonctoriels en $M$, de degré
$0$ pour les structures de $S$-module gradué
\begin{equation}
\label{III.2.1.5.1-fr}
\tag{2.1.5.1}
H^p(U,\widetilde M(*))\xrightarrow{\sim}H^{p+1}((\mathbf f),M)
\qquad\text{pour }p\geq1,
\end{equation}
et une suite exacte fonctorielle en $M$
\begin{equation}
\label{III.2.1.5.2-fr}
\tag{2.1.5.2}
0\longrightarrow H^0((\mathbf f),M)\longrightarrow M
\longrightarrow H^0(U,\widetilde M(*))
\longrightarrow H^1((\mathbf f),M)\longrightarrow0,
\end{equation}
où tous les homomorphismes sont de degré $0$.
\end{proposition}

\begin{proof}
En effet, on a
\[
\Gamma(U_{i_0i_1\ldots i_p},(M(n))^\sim)
=(M_{f_{i_0}\cdots f_{i_p}})_n
\]
par définition (II, 2.5.2), donc
$\Gamma(U_{i_0i_1\ldots i_p},\widetilde M(*))
=M_{f_{i_0}\cdots f_{i_p}}$. Le reste du raisonnement est alors le même
que pour démontrer (\hyperref[III.2.1.4-fr]{2.1.4}), tenant compte de ce
que $X$ est un schéma.
\end{proof}

\begin{env}[2.1.6]
\label{III.2.1.6-fr}
\emph{Remarques.} ---
(i) Sous les conditions de (\hyperref[III.2.1.5-fr]{2.1.5}), les foncteurs
$\Gamma(U_{i_0i_1\ldots i_p},\widetilde M(*))$ sont exacts en $M$, en
vertu de (0\textsubscript{I}, 1.3.2)~; le même raisonnement que dans
(1.2.5) montre alors que si
$0\to M'\to M\to M''\to0$ est une suite exacte de $S$-modules gradués
(où les homomorphismes sont de degré $0$), on a des diagrammes
commutatifs pour tout $p\geq0$
\begin{equation}
\label{III.2.1.6.1-fr}
\tag{2.1.6.1}
\xymatrix{
H^p(U,\widetilde {M''}(*))\ar[r]^-\partial\ar[d] &
H^{p+1}(U,\widetilde {M'}(*))\ar[d]\\
H^{p+1}((\mathbf f),M'')\ar[r]^-\partial &
H^{p+2}((\mathbf f),M').
}
\end{equation}

(ii) La prop. (\hyperref[III.2.1.5-fr]{2.1.5}) sera surtout intéressante
lorsque $S$ sera une $A$-algèbre engendrée par un nombre fini d'éléments de
degré $1$, $A$ étant supposé noethérien~; en
